\documentclass[]{UCD_CS_41820_Report}
\usepackage{graphicx}
\usepackage{listings}
\usepackage{xcolor}
\usepackage{booktabs}
\usepackage{longtable}
\usepackage{titlesec}
\usepackage{amsmath}
\definecolor{codegreen}{rgb}{0,0.6,0}
\definecolor{codegray}{rgb}{0.5,0.5,0.5}
\definecolor{codepurple}{rgb}{0.58,0,0.82}
\definecolor{backcolour}{rgb}{0.95,0.95,0.92}
\usepackage{color}
\definecolor{lightgray}{rgb}{0.95, 0.95, 0.95}
\definecolor{darkgray}{rgb}{0.4, 0.4, 0.4}
\definecolor{editorGray}{rgb}{0.95, 0.95, 0.95}
\definecolor{editorOcher}{rgb}{1, 0.5, 0}
\definecolor{editorGreen}{rgb}{0, 0.5, 0}
\definecolor{orange}{rgb}{1,0.45,0.13}
\definecolor{olive}{rgb}{0.17,0.59,0.20}
\definecolor{brown}{rgb}{0.69,0.31,0.31}
\definecolor{purple}{rgb}{0.38,0.18,0.81}
\definecolor{lightblue}{rgb}{0.1,0.57,0.7}
\definecolor{lightred}{rgb}{1,0.4,0.5}
\usepackage{upquote}
\usepackage{textcomp}
\usepackage{float}
\setcounter{secnumdepth}{3}
\titlespacing*{\section}{0pt}{10pt}{8pt}
\titlespacing*{\subsection}{0pt}{10pt}{4pt}
% CSS
\lstdefinelanguage{CSS}{
  keywords={color,background-image:,margin,padding,font,weight,display,position,
    top,left,right,bottom,list,style,border,size,white,space,min,width,
    transition:, transform:, transition-property, transition-duration,
    transition-timing-function},
  sensitive=true,
  morecomment=[l]{//},
  morecomment=[s]{/*}{*/},
  morestring=[b]',
  morestring=[b]",
  alsoletter={:},
  alsodigit={-}
}

% JavaScript
\lstdefinelanguage{JavaScript}{
  morekeywords={typeof, new, true, false, catch, function, return, null, switch,
    var, let, const, if, in, while, do, else, case, break, await, async},
  morecomment=[s]{/*}{*/},
  morecomment=[l]//,
  morestring=[b]",
  morestring=[b]',
  morestring=[b]`,
  literate={\$}{{$$}}1
}

\lstdefinelanguage{HTML5}{
  language=html,
  sensitive=true,
  alsoletter={<>=-},
  morecomment=[s]{<!-}{-->},
  tag=[s],
  otherkeywords={>,<!DOCTYPE,</html,<html,<head,<title,</title,<style,</style,
    <link,</head,<meta,/>,</body,<body,</div,<div,</div>,</p,<p,</p>,
    </script,<script,<canvas,/canvas>,<svg,<rect,<animateTransform,</rect>,
    </svg>,<video,<source,<iframe,</iframe>,</video>,<image,</image>,
    <header,</header,<article,</article},
  ndkeywords={=,charset=,src=,id=,width=,height=,style=,type=,rel=,href=,
    fill=,attributeName=,begin=,dur=,from=,to=,poster=,controls=,x=,y=,
    repeatCount=,xlink:href=,margin:,padding:,background-image:,border:,
    top:,left:,position:,width:,height:,margin-top:,margin-bottom:,
    font-size:,line-height:,transform:,-moz-transform:,-webkit-transform:,
    animation:,-webkit-animation:,transition:,transition-duration:,
    transition-property:,transition-timing-function:}
}

\lstdefinestyle{htmlcssjs}{%
  basicstyle={\footnotesize\ttfamily},
  frame=b,
  xleftmargin={0.75cm},
  numbers=left,
  stepnumber=1,
  firstnumber=1,
  numberfirstline=true,
  identifierstyle=\color{black},
  keywordstyle=\color{blue}\bfseries,
  ndkeywordstyle=\color{editorGreen}\bfseries,
  stringstyle=\color{editorOcher}\ttfamily,
  commentstyle=\color{brown}\ttfamily,
  language=HTML5,
  alsolanguage=JavaScript,
  alsodigit={.:;},
  tabsize=2,
  showtabs=false,
  showspaces=false,
  showstringspaces=false,
  extendedchars=true,
  breaklines=true,
}

\lstdefinestyle{py}{%
  language=python,
  basicstyle=\footnotesize\ttfamily,
  numbers=left,
  numbersep=5pt,
  tabsize=4,
  extendedchars=true,
  breaklines=true,
  keywordstyle=\color{blue}\bfseries,
  frame=b,
  commentstyle=\color{brown}\itshape,
  stringstyle=\color{editorOcher}\ttfamily,
  showspaces=false,
  showtabs=false,
  xleftmargin=17pt,
  framexleftmargin=17pt,
  framexrightmargin=5pt,
  framexbottommargin=4pt,
  showstringspaces=false,
}

\lstdefinestyle{mystyle}{
  backgroundcolor=\color{backcolour},
  commentstyle=\color{codegreen},
  keywordstyle=\color{magenta},
  numberstyle=\tiny\color{codegray},
  stringstyle=\color{codepurple},
  basicstyle=\ttfamily\footnotesize,
  breakatwhitespace=false,
  breaklines=true,
  captionpos=b,
  keepspaces=true,
  numbers=left,
  numbersep=5pt,
  showspaces=false,
  showstringspaces=false,
  showtabs=false,
  tabsize=2
}

\lstset{style=htmlcssjs}
\usepackage{hyperref}
\hypersetup{
  colorlinks=true,
  linkcolor=blue,
  filecolor=magenta,
  urlcolor=cyan,
}
\lstset{mathescape=true, basicstyle=\ttfamily}

\usepackage[backend=biber,style=nature]{biblatex}
\addbibresource{references.bib}

%%% Project details
\def\firststudentname{Ashwin Prasanth}
\def\firststudentid{25234753}
\def\secondstudentname{Kiran Meenakshi Sundaram}
\def\secondstudentid{25217694}
\def\projecttitle{Ethical Analysis of the Epic Sepsis Model}
\def\supervisorname{Vivek Nallur}

%=======================================================================
\begin{document}
\maketitle
\tableofcontents

%=======================================================================
\chapter{Introduction}
%=======================================================================

Sepsis is a life-threatening condition caused by a dysregulated host response to
infection and is associated with high mortality and substantial healthcare costs.
In hospitalized patients, sepsis occurs in approximately 2--6\% of admissions, with
reported mortality rates ranging from 15\% to 40\% depending on severity and the
time of intervention.

Clinical evidence consistently shows that earlier identification and treatment, particularly
timely administration of antibiotics, significantly reduce mortality. As a result,
healthcare systems have increasingly adopted AI-enabled Clinical Decision Support
Systems (CDSS) to assist clinicians in early detection and risk stratification of sepsis.

This project focuses on a real-world, deployed AI-enabled CDSS: the \textbf{Epic Sepsis
Model (ESM)} Inpatient Predictive Analytic Tool, integrated into the Epic electronic
health record (EHR) system. Epic Systems is the largest EHR provider in the United
States, serving approximately 54\% of patients nationwide, making the ESM one of the
most widely deployed AI-based decision support tools in clinical practice.

The ESM is designed to continuously monitor hospitalized patients and generate a
numerical risk score indicating the likelihood of sepsis, enabling earlier clinical
intervention.

Unlike experimental or research-only models, the ESM operates in live hospital
environments and directly influences clinical workflows. Its outputs trigger alerts
to nurses and physicians when a predefined risk threshold is exceeded, thereby
shaping diagnostic attention, treatment prioritisation, and escalation decisions.
The ESM therefore represents a \emph{safety-critical} AI system whose ethical
implications extend beyond model accuracy to issues of trust, accountability,
over-reliance, and human--AI collaboration.
\vspace{-8pt}
%-----------------------------------------------------------------------
\subsection*{Key Features of the System}
%-----------------------------------------------------------------------

The Epic Sepsis Model is a data-driven, sub-symbolic AI system implemented within
a commercial electronic health record platform. Published validation studies provide
sufficient evidence to characterise its technical structure and behaviour.

\subsubsection*{AI Technology and Model Characteristics}

\begin{table}[H]
\centering
\renewcommand{\arraystretch}{1.3}
\begin{tabular}{@{}ll@{}}
\toprule
\textbf{Property} & \textbf{Detail} \\
\midrule
Model type      & Data-driven predictive model (sub-symbolic) \\
Input features  & $\approx$80 clinical and demographic variables \\
                & (vital signs, lab results, comorbidities, demographics) \\
Training data   & $\approx$500{,}000 patient encounters \\
Output          & Continuously updated numerical sepsis risk score \\
Deployment      & Fully integrated into Epic EHR; operates in real time \\
\bottomrule
\end{tabular}
\caption{Summary of Epic Sepsis Model technical characteristics.}
\label{tab:esm_features}
\end{table}

The ESM does not rely on explicit clinical rules, symbolic reasoning, or
hand-crafted knowledge representations. Patient data are combined and weighted
according to learned statistical relationships to produce a risk score. While the
system can surface which variables contributed most to an elevated score, its
internal decision logic is not fully transparent to end users, placing it firmly
in the category of \emph{sub-symbolic, black-box} clinical decision support.
\setlength{\parskip}{0pt}
%-----------------------------------------------------------------------
\subsection*{Operational Threshold and Performance}
%-----------------------------------------------------------------------

In the validation study conducted at a 746-bed academic medical centre, an alert
threshold of \textbf{ESM\,$\geq$\,5} was selected using receiver operating characteristic
(ROC) curve analysis. Table~\ref{tab:esm_perf} summarises the performance at this
threshold.

\begin{table}[H]
\centering
\renewcommand{\arraystretch}{1.3}
\begin{tabular}{@{}lc@{}}
\toprule
\textbf{Metric} & \textbf{Value} \\
\midrule
Area Under ROC Curve (AUC)      & 0.834 \\
Sensitivity                     & 86.0\% \\
Specificity                     & 80.8\% \\
Positive Predictive Value (PPV) & 33.8\% \\
Negative Predictive Value (NPV) & 98.1\% \\
\bottomrule
\end{tabular}
\caption{ESM performance at alert threshold $\geq 5$.}
\label{tab:esm_perf}
\end{table}

These results indicate that the ESM is optimised as a \emph{screening} tool,
prioritising sensitivity and early detection rather than definitive diagnosis.
\setlength{\parskip}{0pt}
%=======================================================================
\section{Functional Goals of the System}
%=======================================================================

The primary functional goal of the Epic Sepsis Model is the early identification of
hospitalized patients at risk of developing sepsis, enabling clinicians to initiate
timely evaluation and treatment. This goal is operationalised through continuous risk
scoring and automated alerting within the EHR.
\setlength{\parskip}{0pt}
\paragraph{1.\ Detect elevated sepsis risk earlier than routine clinical recognition.}
By continuously analysing patient data, the ESM aims to identify deterioration before
overt clinical signs prompt clinician action.

\paragraph{2.\ Prompt timely clinical intervention.}
When the ESM score exceeds the alert threshold, clinicians are notified and expected
to assess the patient, consider sepsis protocols, and initiate appropriate management.

\paragraph{3.\ Reduce time to antibiotic administration.}
In the post-implementation phase of the validation study, the median time from alert
to antibiotic administration decreased from 150 minutes to 90 minutes, and the
proportion of patients receiving antibiotics within three hours increased from 55.6\%
to 69.8\%.

\paragraph{4.\ Improve patient outcomes, particularly mortality.}
Among patients with ESM scores $\geq 5$ who had not already received antibiotics,
implementation was associated with a \textbf{44\% reduction} in the odds of
in-hospital sepsis-related mortality (odds ratio 0.56; 95\% CI 0.39--0.80).

\paragraph{5.\ Integrate seamlessly into existing clinical workflows.}
The system is embedded within the Epic EHR and requires no separate interface,
enabling widespread deployment with minimal disruption.

%-----------------------------------------------------------------------
\subsection*{Ethical Relevance of the Functional Design}
%-----------------------------------------------------------------------

Although the ESM demonstrably improves early intervention and is associated with
reduced mortality, its design raises ethical challenges beyond functional performance.
The system produces risk scores that significantly influence clinician behaviour,
yet its internal reasoning is largely opaque. Clinicians remain legally and morally
responsible for all decisions, while the AI system shapes attention, urgency, and
treatment timing. This creates potential risks of \emph{automation bias},
\emph{over-reliance}, and \emph{accountability gaps} in high-stakes clinical contexts.

These considerations motivate the focus of this project: not on building a new
predictive model, but on explicitly identifying ethical goals, formalising them
as constraints or checks, and demonstrating---through targeted scripts and validation
cases---when such goals may or may not be satisfied.

%=======================================================================
\chapter{Stakeholder Identification}
%=======================================================================

The Epic Sepsis Model (ESM) operates within a complex hospital environment involving
clinical personnel, patients, healthcare institutions, and technology providers.
As a deployed AI-enabled Clinical Decision Support System (CDSS), its ethical impact
arises from interactions between these stakeholders rather than from the predictive
model alone.

This chapter identifies the primary and secondary stakeholders affected by the system
and maps them across social, technical, economic, environmental, political and legal
dimensions.

\vspace{-8pt}
\section{Stakeholders of the Epic Sepsis Model}

\subsection*{Primary Stakeholders}

Primary stakeholders are those who directly interact with the system or are directly
affected by its outputs.

\subsubsection*{1.\ Clinicians (Physicians and Nurses)}
\textit{Category: Social, Technical, Legal}

Clinicians are the primary users of the ESM. The system generates alerts when a
patient's sepsis risk score exceeds a predefined threshold (\textbf{ESM\,$\geq$\,5}),
prompting clinical evaluation and intervention.

Although clinicians retain final decision-making authority, the alerts influence
attention, prioritisation, and timing of treatment. Clinicians are also legally
accountable for patient outcomes, even when decisions are informed by AI-generated
risk scores.

\subsubsection*{2.\ Patients}
\textit{Category: Social, Ethical, Legal}

Hospitalised patients are the subjects of the predictive model. They do not directly
interact with the system but are affected by its outputs through changes in clinical
behaviour, including earlier antibiotic administration and ICU escalation.

Patients bear the risks of false positives (unnecessary treatment, antibiotic overuse)
and false negatives (missed or delayed diagnosis).

\subsubsection*{3.\ Hospital and Healthcare Providers}
\textit{Category: Economic, Technical, Legal}

Hospitals deploy the ESM as part of the Epic EHR infrastructure. They benefit from
potential reductions in mortality, length of stay, and compliance with sepsis quality
measures. At the same time, hospitals assume institutional liability and must manage
alert fatigue, workflow disruption, and training requirements.

\subsection*{Secondary Stakeholders}

Secondary stakeholders influence the system indirectly or are affected at a broader
systemic level.

\subsubsection*{4.\ Epic Systems Corporation}
\textit{Category: Technical, Economic, Legal}

Epic develops and maintains the proprietary ESM algorithm. Its design choices
including feature selection, model updates and alert presentation shape how clinicians
interpret and act on risk scores.

\subsubsection*{5.\ Health Regulators and Accreditation Bodies}
\textit{Category: Political, Legal}

Regulatory bodies influence adoption of sepsis detection systems through quality
standards, reporting requirements, and patient safety regulations. Although the ESM
is not publicly regulated as a medical device, its use affects compliance with
national sepsis guidelines.

\subsubsection*{6.\ Society and Public Health Systems}
\textit{Category: Social, Economic}

At a population level, widespread deployment of sepsis detection tools influences
healthcare costs, antibiotic stewardship, and public trust in AI-assisted medicine.

\subsubsection*{7.\ Environmental and Infrastructure Considerations}
\textit{Category: Environmental, Technical}

Large scale deployment across hospital networks contributes to cumulative energy
consumption and infrastructure expansion. While individual risk score computations
have negligible marginal environmental impact, responsible digital health governance
must consider long-term sustainability of computational infrastructure.

%=======================================================================
\section{Ethical Goals of the System}
%=======================================================================

The Epic Sepsis Model does not publish an explicit ethical framework or set of
ethical principles. Public documentation and validation studies focus primarily
on clinical effectiveness, workflow integration, and patient outcomes.

There are no explicit public statements addressing transparency, accountability,
or ethical safeguards beyond standard clinical governance practices.

This absence is ethically relevant, as it requires ethical safeguards to be inferred
from system behaviour and deployment structure rather than from explicitly articulated
normative commitments.

%=======================================================================
\section{Ethical Considerations and Schools of Thought}
%=======================================================================

For this project, three ethical frameworks are used to analyse the system:

\begin{itemize}
  \item \textbf{Utilitarianism}
  \item \textbf{Deontological Ethics}
  \item \textbf{Virtue Ethics}
\end{itemize}

These were selected because they capture outcome-based, duty-based, and professional
character-based ethical concerns relevant to clinical AI systems.

\subsection*{1.\ Utilitarian Ethics (Outcome-Oriented)}

\textbf{Core Principle:} Actions are ethically justified if they maximise overall
benefit and minimise harm across all affected parties (Bentham, Mill).

\paragraph{Relevant Variables and System Functionality.}
Sepsis mortality rates across the patient population; time to antibiotic
administration per alert; false positive rate (FPR) and its burden on clinical
workflow; hospital-wide resource utilisation.

\paragraph{Ethical Goal U1: Minimise Preventable Sepsis Mortality.}~\\

\textbf{Actionable Principle (Utilitarian: Greatest Aggregate Benefit):}
The ESM alert system is ethically justified only if the reduction in
sepsis-related mortality across the treated population outweighs the aggregate
harm caused by false positive alerts, including unnecessary antibiotic
administration, increased costs, and alert fatigue-induced inattention. Vasey
et al.\ \cite{vasey2022} (DECIDE-AI) demonstrate that ethical risks in AI-based
CDSS frequently arise from deployment failures rather than model errors, meaning
the utilitarian calculation must be evaluated continuously in real-world conditions,
not only at the point of validation.\par

\medskip
\textbf{Measurable Condition:} The system satisfies this principle if the odds
ratio of sepsis-related mortality in the alert-active period remains below 1.0 at
a statistically significant level ($p < 0.05$). The published validation study
\cite{cull2023} reports OR\,=\,0.56 (95\% CI 0.39--0.80), satisfying this condition
at the studied institution.\par

\medskip
\textbf{Violation Trigger:} If post-deployment monitoring reveals that the false
positive burden has increased clinician override rates to the point where true
positive alerts are being systematically ignored, the net utility calculation is
violated. Specifically, if the proportion of true positive alerts acted upon falls
below the proportion in the pre-implementation period, the system no longer produces
net benefit and the utilitarian justification for its deployment fails.\par

\paragraph{Ethical Goal U2: Control False Positive Burden.}~\\

\textbf{Actionable Principle (Utilitarian : Harm Minimisation):}
The false positive rate must be controlled such that the cumulative harm of
unnecessary interventions does not outweigh the benefit of early detection. A system
with sensitivity optimised at 86.0\% and PPV of 33.8\% accepts that approximately
two in three alerts are false positives \cite{cull2023}. This trade-off is ethically
acceptable only if alert fatigue does not reduce the clinical response rate to true
positive alerts. Shwedeh and Alzoubi \cite{shwedeh2025} demonstrate empirically that
over-reliance and alert fatigue emerge specifically in high-pressure deployment
environments where governance and oversight are insufficient, reinforcing that this
condition must be actively monitored rather than assumed.\par

\medskip
\textbf{Measurable Condition:} Alert fatigue is operationally defined as a clinician
override rate exceeding 70\% across consecutive alert sessions. If this threshold is
crossed, the PPV of acted-upon alerts falls below the level required for net
population benefit.\par

\medskip
\textbf{Violation Trigger:} Override rate $>$ 70\% across a rolling window of alerts,
or a statistically significant increase in time-to-antibiotic following the
introduction of the alert system.\par

\subsection*{2.\ Deontological Ethics (Duty and Responsibility)}

\textbf{Core Principle:} Ethical action is determined by adherence to duties and
rules, independent of outcomes. Agents must be treated as rational ends in
themselves, not merely as means (Kant, \textit{Groundwork of the Metaphysics of
Morals}).

\paragraph{Relevant Variables and System Functionality.}
Availability of model reasoning to clinicians; clinician legal and moral
accountability for alert-informed decisions; governance requirement for alert
acknowledgement.

\paragraph{Ethical Goal D1: Preserve Clinician Epistemic Access (Duty of Non-Deception).}~\\

\textbf{Actionable Principle (Kantian : Respect for Rational Agency):}
A clinician who is held legally and morally responsible for a clinical decision must
have meaningful epistemic access to the basis of the information that prompted that
decision. Denying access to model reasoning while retaining accountability violates
the Kantian duty to treat rational agents as ends in themselves,  the clinician is
reduced to a conduit for automated output rather than a deliberating professional.
Amann et al.\ \cite{amann2022} establish that lack of explainability causes systems
nominally classified as decision support to drift toward decision determining,
undermining meaningful human oversight. Xu et al.\ \cite{xu2023} further argue that
interpretability is a fundamental ethical requirement for AI-based CDSS, not an
optional feature, and that post-hoc explanation methods alone are insufficient to
guarantee accountability. \v{C}artolovni et al.\ \cite{cartoloni2022} identify the
resulting accountability gap where formal responsibility remains with the clinician
while practical epistemic control has shifted to the AI as the dominant ethical
concern in the AI-based medical decision support literature.\par

\medskip
\textbf{Measurable Condition:} The system satisfies this principle if, for every
alert fired, the clinician can access at minimum the top contributing feature
variables and their directional influence on the score. The current ESM does not
satisfy this condition, feature-level reasoning is not exposed to bedside
clinicians.\par

\medskip
\textbf{Violation Trigger:} Any alert fired without accompanying feature-level
explanation constitutes a violation of this principle. In the current deployment,
this condition is violated for 100\% of alerts.\par

\paragraph{Ethical Goal D2: Enforce Governance Accountability (Duty of Institutional Responsibility).}~\\

\textbf{Actionable Principle (Deontological : Rule-Based Governance):}
Every alert fired by the system must be formally acknowledged by a responsible
clinician. An unacknowledged alert represents a failure of the institutional duty
of care, the system has identified a patient at risk and no accountable agent has
accepted responsibility for the assessment.\par

\medskip
\textbf{Measurable Condition:} Governance is satisfied if and only if for all alerts
$A$ fired by the system, there exists a corresponding clinician acknowledgement $C$
within a defined response window $T$. Formally: $Alert(A) \rightarrow
Acknowledged(C, T)$.\par

\medskip
\textbf{Violation Trigger:} Any alert that passes without documented clinician
acknowledgement within the response window violates this rule. This is
operationalised as the \texttt{governance\_check} constraint in the implementation
and verified formally using Z3 SMT.\par

\subsection*{3.\ Virtue Ethics (Professional Character and Practice)}

\textbf{Core Principle:} Ethical systems should cultivate and support good
professional character such as practical wisdom (\textit{phronesis}),
attentiveness, and calibrated judgement rather than replace them (Aristotle,
\textit{Nicomachean Ethics}).

\paragraph{Relevant Variables and System Functionality.}
Clinician override and compliance patterns over time; alert frequency and its
effect on attentiveness; long-term effect of AI-assisted workflows on clinical skill.

\paragraph{Ethical Goal V1: Prevent Automation Bias and Preserve Clinical Prudence.}~\\

\textbf{Actionable Principle (Virtue Ethics : Phronesis):}
A clinical AI system must be designed and deployed in a way that supports rather
than displaces the exercise of practical wisdom. Automation bias: the tendency
to follow AI output without independent assessment directly undermines
\textit{phronesis}. A system that produces this effect at scale is ethically
deficient regardless of its population-level mortality benefit. Panigutti et al.\
\cite{panigutti2022} provide empirical evidence that AI explanations systematically
increase human reliance on AI advice regardless of whether that reliance improves
decision accuracy, demonstrating that the problem of automation bias cannot be
resolved by adding explanations alone. \v{C}artolovni et al.\ \cite{cartoloni2022}
identify automation bias and deskilling as the most clinician-specific ethical risks
in AI-based medical decision support, directly implicating virtue ethics as the
relevant framework.\par

\medskip
\textbf{Measurable Condition:} Automation bias is operationally present if the
clinician compliance rate with alerts defined as the proportion of alerts followed
without documented independent clinical assessment exceeds 90\% across a rolling
observation window.\par

\medskip
\textbf{Violation Trigger:} Compliance rate $>$ 90\% without documented independent
assessment, operationalised as the \texttt{detect\_over\_reliance} constraint in the
implementation (threshold: \texttt{compliance\_rate > 0.9}).\par

\paragraph{Ethical Goal V2: Prevent Clinical Deskilling.}~\\

\textbf{Actionable Principle (Virtue Ethics : Maintenance of Professional Competence):}
Repeated reliance on AI-generated alerts without active clinical reasoning
constitutes a long-term threat to the professional competence required for safe
independent practice. Clinicians who have not practiced independent sepsis assessment
over an extended period may lose the diagnostic calibration that makes meaningful
human oversight possible. \v{C}artolovni et al.\ \cite{cartoloni2022} explicitly
identify deskilling and loss of epistemic authority as physician-specific ethical
risks in AI-based decision support, distinct from automation bias a clinician may
not be biased in any given moment but may have cumulatively lost the skill required
for genuine independent assessment.\par

\medskip
\textbf{Measurable Condition:} Deskilling risk is flagged if clinician independent
override decisions, cases where a clinician actively disagrees with and documents
a reasoned rejection of an alert fall below a defined minimum frequency per
clinician per month.\par

\medskip
\textbf{Violation Trigger:} Independent override rate $<$ minimum threshold per
clinician per defined period, indicating habitual passive compliance rather than
active clinical engagement.\par

\subsection*{Ethical Framework--Stakeholder Alignment}

The three ethical schools correspond to different stakeholder perspectives.
Utilitarian analysis primarily concerns patient outcomes and hospital-level mortality
reduction. Deontological considerations are most salient for clinicians, who retain
formal decision authority and legal responsibility for AI-informed actions. Virtue
ethics focuses on the professional character of clinicians interacting repeatedly
with alert systems, emphasising calibration, attentiveness, and avoidance of
automation bias. Together, these frameworks provide a pluralistic ethical lens
appropriate for socio-technical clinical systems.

%=======================================================================
\section{Values}
%=======================================================================

\subsection*{Stakeholder Values and Alignment}

\begin{table}[H]
\centering
\renewcommand{\arraystretch}{1.3}
\begin{tabular}{@{}p{3cm}p{4cm}p{3cm}p{4cm}@{}}
\toprule
\textbf{Stakeholder} & \textbf{Primary Values} & \textbf{Alignment with Ethical Goals} & \textbf{Conflict} \\
\midrule
Clinicians
  & Professional autonomy, accountability, patient safety
  & D1, D2, V1, V2
  & Conflicts with Epic's proprietary opacity (D1) \\[4pt]
Patients
  & Timely treatment, equitable care, safety
  & U1, U2, D1
  & No direct conflict benefit from all goals \\[4pt]
Hospitals
  & Efficiency, liability reduction, compliance
  & U1, U2, D2
  & Cost of explainability features conflicts with U2 \\[4pt]
Epic Systems
  & Proprietary protection, commercial viability
  & None stated publicly
  & Directly conflicts with D1 (transparency) \\[4pt]
Regulators
  & Safety, standardisation, accountability
  & D2, U1
  & May conflict with speed of deployment \\[4pt]
Society
  & Equitable access, antibiotic stewardship
  & U2, V1
  & Alert fatigue burden conflicts with U2 \\
\bottomrule
\end{tabular}
\caption{Stakeholder values, ethical goal alignment, and conflicts.}
\label{tab:stakeholder_values}
\end{table}

\subsection*{Values Requiring Negotiation}

The most significant value conflict is between Epic's commercial interest in
maintaining proprietary opacity and the deontological requirement for clinician
epistemic access (D1). These cannot be simultaneously satisfied in the current
deployment. Resolving this conflict requires either mandatory explainability
regulation or contractual transparency obligations in hospital procurement.

A secondary conflict exists between the utilitarian justification for high
sensitivity (U1) and the harm minimisation requirement (U2). Maximising sensitivity
necessarily increases false positive burden. This trade-off cannot be resolved by
system design alone, it requires institutional threshold governance involving
clinicians, patients and ethics review.

\subsection*{Prioritisation}

Patient safety values take precedence, as patients are the population most directly
harmed by system failures and least able to advocate within the deployment structure.
Clinician professional values are second, as they are the responsible agents whose
judgement the system most directly affects. Institutional and commercial values,
while legitimate, must not override the first two categories.

\subsection*{Functional-to-Ethical Goal Balance}

The primary functional purpose of the Epic Sepsis Model is early detection and
improved outcomes in sepsis care. Ethical goals constrain how that purpose is
achieved rather than overriding it. For primary stakeholders, functional goals
significantly outnumber ethical constraints, reflecting the system's clinical purpose
while ensuring ethical safeguards remain enforceable.

\begin{table}[H]
\centering
\renewcommand{\arraystretch}{1.3}
\begin{tabular}{@{}ll@{}}
\toprule
\textbf{Functional Goals} & \textbf{Ethical Goals} \\
\midrule
Early detection                      & Prevent over-reliance (V1) \\
Threshold-based alerting             & Preserve epistemic access (D1) \\
Reduced mortality                    & Enforce governance accountability (D2) \\
Workflow integration                 & Control false positive burden (U2) \\
Antibiotic timing improvement        & Prevent clinical deskilling (V2) \\
\bottomrule
\end{tabular}
\caption{Functional vs.\ ethical goals of the ESM.}
\label{tab:goal_balance}
\end{table}

Ethical goals are more difficult to implement and validate but do not overshadow
the system's clinical objectives.
%=======================================================================
\chapter{Implementation}
%=======================================================================

\textbf{Link to Repository:}
\href{https://github.com/AshwinPrasanth/AI--Ethics}{https://github.com/AshwinPrasanth/AI--Ethics}

The repository contains a structured simulation and formal verification framework
modelling the decision-structural properties of the Epic Sepsis Model. The objective
is not architectural replication of the proprietary system, but formalisation and
validation of functional and ethical properties. The implementation consists of two
validation layers: a \textbf{behavioural simulation layer} and a \textbf{formal
logical verification layer} (Z3 SMT). Together, they operationalise ethical
evaluation as executable constraint verification.

%=======================================================================
\section{Description of Approach}
%=======================================================================

\subsection{Model Abstraction}

The Epic Sepsis Model is a proprietary, data-driven predictive system trained on
approximately 500,000 patient encounters and incorporating $\sim$80 clinical
variables. Its internal architecture, feature weighting, and training pipeline
are not publicly disclosed.

This implementation does not attempt model reproduction. Instead, it preserves the
operational invariants that define the system's decision structure: continuous
numerical risk scoring; threshold-based alert triggering (Score\,$\geq$\,5);
screening-oriented behaviour (high sensitivity design); non-zero false positive
possibility; human-in-the-loop decision authority; and integration into clinician
workflow. The system therefore represents a structural abstraction of a sub-symbolic
CDSS, sufficient for ethical validation experiments.

%-----------------------------------------------------------------------
\subsection{Domain Representation}

Each synthetic inpatient instance is represented as:
\[
  x = (T,\; HR,\; Lactate,\; WBC,\; SBP,\; PriorAntibiotics)
\]

where $T$ is temperature, $HR$ heart rate, $WBC$ white blood cell count, and $SBP$
systolic blood pressure. These represent a clinically meaningful subset of ESM inputs.
The risk score is computed as:
\[
  Score = \sum_i w_i x_i
\]
with fixed weights $w_i$. The purpose is not predictive realism, but simulation of
sub-symbolic weighted aggregation of heterogeneous physiological signals.

%-----------------------------------------------------------------------
\subsection{Functional Goal Encoding}

Functional goals are encoded in \texttt{functional\_validation.py}.

\paragraph{F1: Continuous Risk Stratification.}
The system produces a real-valued risk score for each patient. High-risk synthetic
cases yield $Score \geq 5$; moderate-risk cases yield $Score < 5$.

\paragraph{F2: Deterministic Alert Trigger.}
$Alert = (Score \geq 5)$, aligning with the published ESM cut-point.

\paragraph{F3: Workflow Activation.}
If $Alert = \text{True}$, the system simulates escalation to a clinician and a
decision node requiring response, modelling the real-world operational role of
ESM alerts.

%-----------------------------------------------------------------------
\subsection{Ethical Goal Formalisation (Simulation Layer)}

Ethical goals are encoded as executable constraint checks in
\texttt{ethical\_validation.py}.

\paragraph{E1: False Positive Monitoring (Utilitarian Trade-Off).}
$Alert = \text{True} \land GroundTruth = \text{False}$. The False Positive Rate is
computed as $FPR = FP / (FP + TN)$. If $FPR > 0$, alert fatigue risk is flagged
and logged.

\paragraph{E2: Over-Reliance Detection (Automation Bias).}
$Alert = \text{True} \land ClinicianAction = Accept \land GroundTruth = \text{False}$.
Captures blind compliance with AI output.

\paragraph{E3: Under-Reliance Detection.}
$Alert = \text{True} \land \mathit{ClinicianAction} = \mathit{Ignore}$ triggers\\
 $\mathit{EthicalGoalSatisfied} = \text{False}$.

\paragraph{E4: Responsibility Gap Detection.}
$Alert = \text{True} \land ExplanationUnavailable$. \textit{Known limitation:}
In the implementation, \texttt{generate\_explanation()} always returns a feature
contribution dictionary, so the gap is never triggered in simulation. The real ESM
does not expose its internal reasoning; the Z3 layer addresses this structurally.

\paragraph{E5: Alert Fatigue Monitoring.}
Repeated alerts accumulate across a session. If counter $> 5$, fatigue risk is
flagged.

\paragraph{E6: Governance Constraint.}
$Alert \rightarrow Acknowledged$. Violation triggers a governance failure report.

\paragraph{E7: Data Quality Failure.}
Corrupted or incomplete patient inputs trigger exceptions and logging. When
\texttt{temperature: None} is passed, a \texttt{ValueError} is raised and written
to \texttt{audit.log}.

\paragraph{E8: Age-Based Fairness Scenario.}
Two clinically equivalent patients differing only in age category are evaluated.
If $EquivalentClinicalState \rightarrow (Alert_{elderly} = Alert_{young})$ is
violated, subgroup disparity is flagged. Because the current scoring function does
not incorporate \texttt{age\_group} as a feature, the simulation always returns
equivalent scores; the Z3 layer handles the formal violation case.

%-----------------------------------------------------------------------
\subsection{Formal Logical Verification (Z3 SMT Layer)}
Selected ethical constraints are encoded as logical formulas and analysed using the
Z3 SMT solver in \texttt{z3\_ethics\_verification.py}. Constraints are registered
with \texttt{assert\_and\_track}, enabling UNSAT core extraction.

\paragraph{SAT Case: Duty-of-Care Consistency.}
$(Alert \land GroundTruth) \rightarrow ClinicianResponded$. Z3 returns SAT,
confirming logical consistency.

\paragraph{UNSAT Case: Capacity Conflict.}
Two ethical mandates ($Treat_A \land Treat_B$) conflict with the physical constraint
$\neg(Treat_A \land Treat_B)$. Z3 returns UNSAT and extracts the minimal core:
\{\texttt{Ethical\_Treat\_A}, \texttt{Ethical\_Treat\_B}, \texttt{Capacity\_Constraint}\}.

\paragraph{UNSAT Case: Age-Based Fairness Violation.}
$EquivalentClinicalState \rightarrow (Alert_{elderly} = Alert_{young})$. Violation
produces UNSAT; core isolates \{\texttt{Elderly\_Alert\_True},
\texttt{Equivalent\_Group\_True}, \texttt{Fairness\_Symmetry},
\texttt{Young\_Alert\_False}\}.

\paragraph{UNSAT Case: Governance Violation.}
$Alert \rightarrow Acknowledged$. UNSAT core: \{\texttt{Governance\_Rule},
\texttt{Alert\_Occurred}, \texttt{No\_Acknowledgement}\}.

\paragraph{Architectural significance.}
The SMT layer detects structural ethical contradictions, extracts minimal conflicting
constraint sets, distinguishes behavioural violation from logical impossibility,
and complements simulation-based evaluation.

%-----------------------------------------------------------------------
\subsection{Scope Delimitation}

The implementation does not replicate proprietary ESM weighting, perform ROC
optimisation, estimate AUC, model mortality causality, or use temporal logic.
The system focuses on threshold-based decision structure and constraint consistency.

%=======================================================================
\section{Points of Conflict}
%=======================================================================

\subsection*{Conflict 1: Fairness vs.\ Clinical Equivalence}

\textbf{Original decision:} The initial scoring model treated all patients
identically using only physiological variables, with no subgroup differentiation.\par
\medskip
\textbf{Value conflict:} Deontological and virtue ethics considerations raised
the question of whether a system that does not explicitly test for subgroup equity
could produce discriminatory outcomes for elderly patients with atypical sepsis
physiology.\par
\medskip

\textbf{Design change:} Two synthetic patient scenarios namely, \texttt{elderly\_high\_risk\_patient()}
and \\ \texttt{young\_high\_risk\_patient()} were introduced. A fairness check (E8)
was added to the simulation layer and a formal fairness symmetry constraint was
encoded in the Z3 layer.\par
\medskip

\textbf{Stakeholders who benefit:} Patients (particularly elderly patients);
regulators requiring evidence of non-discriminatory behaviour.\par
\medskip

\textbf{Residual limitation:} The current scoring function does not use
\texttt{age\_group} as a feature, so the simulation always returns equivalent
outcomes. The Z3 layer formally proves what would happen if unequal outcomes
occurred.

\subsection*{Conflict 2: Transparency vs.\ Proprietary Opacity}

\textbf{Original decision:} The abstraction included a \texttt{generate\_explanation()}
function returning feature contributions, treating interpretability as available.\par
\medskip

\textbf{Value conflict:} This misrepresented the real ESM. Deontological concerns
around accountability required the abstraction to reflect opacity rather than paper
over it.\par
\medskip

\textbf{Design change:} The E4 constraint was retained but documented as a known
abstraction limitation. The Z3 layer was extended to formally model governance and
accountability constraints that the simulation layer could not fully instantiate.\par
\medskip

\textbf{Stakeholders who benefit:} Clinicians; patients whose safety depends on
clinicians understanding the limits of AI-generated scores.

\subsection*{Conflict 3: Sensitivity Optimisation vs.\ Alert Fatigue}

\textbf{Original decision:} The alert threshold was fixed at Score\,$\geq$\,5
without analysis of how threshold variation affects ethical outcomes.\par
\medskip

\textbf{Value conflict:} Utilitarian analysis identified that optimising for
sensitivity necessarily increases false positives and accumulates alert burden
on clinicians, creating a conflict between population-level mortality benefit and
individual workflow sustainability.\par
\medskip

\textbf{Design change:} A threshold trade-off simulation was added, evaluating
alert behaviour at thresholds of 4, 5, and 7. An alert fatigue counter (E5) and
expanded audit logging were introduced.\par
\medskip

\textbf{Stakeholders who benefit:} Clinicians; hospitals using threshold
configuration as a governance lever; patients, indirectly.

%=======================================================================
\section{Evidence}
%=======================================================================

\subsection*{Functional Validation Output}

Execution of \texttt{python functional\_validation.py} produces:

\noindent\textit{[Insert screenshot: functional\_validation.py terminal output]}

A high-risk profile scores 5.75, correctly triggering an alert. A moderate-risk
profile scores 2.98, correctly suppressing it. This validates functional goals
\textbf{F1} and \textbf{F2}.

%-----------------------------------------------------------------------
\subsection*{Ethical Validation Output}

Execution of \texttt{python ethical\_validation.py} produces:

\noindent\textit{[Insert screenshot: ethical\_validation.py terminal output]}

\begin{table}[H]
\centering
\renewcommand{\arraystretch}{1.3}
\begin{tabular}{@{}lll@{}}
\toprule
\textbf{Output} & \textbf{Ethical Goal} & \textbf{Result} \\
\midrule
Over-reliance detected: True    & E2 --- Automation Bias        & Violation detected \\
Basic FPR: 0.0                  & E1 --- FP Monitoring          & Correct (below threshold) \\
Aggressive FPR: 1.0             & E1 --- FP Monitoring          & Violation flagged \\
Ethical goal satisfied: False   & E3 --- Under-reliance         & Violation detected \\
Responsibility gap: False       & E4 --- Responsibility Gap     & Abstraction limit \\
Alert fatigue: True             & E5 --- Alert Fatigue          & Violation flagged \\
Governance satisfied: False     & E6 --- Governance             & Violation detected \\
Data quality failure            & E7 --- Data Quality           & Exception caught \\
No subgroup disparity           & E8 --- Fairness               & Z3 handles formal case \\
\bottomrule
\end{tabular}
\caption{Ethical validation results mapped to ethical goals.}
\label{tab:ethical_results}
\end{table}

%-----------------------------------------------------------------------
\subsection*{Z3 Formal Verification Output}

Execution of \texttt{python z3\_ethics\_verification.py} produces:

\noindent\textit{[Insert screenshot: z3\_ethics\_verification.py terminal output]}

The SAT result confirms the ethical rule set is consistent under compliant
behaviour. The three UNSAT results prove that fairness violations, triage deadlock,
and governance failures are logically impossible to reconcile with their respective
constraints. UNSAT core extraction identifies the minimal responsible constraint
sets.


%-----------------------------------------------------------------------
\subsection*{Audit Log}

\noindent\textit{[Insert screenshot: audit.log file content]}

The audit log provides a persistent, timestamped accountability trace across
multiple runs. Repeated entries confirm that synthetic scenarios are deterministic
by design. The log records over-reliance warnings, false positive events,
under-reliance violations, and data quality failures, modelling the audit trail
a real deployed CDSS would be expected to maintain for governance purposes.

%=======================================================================
\chapter{Claude Says\ldots}
%=======================================================================

\section*{LLM Details}

\textbf{Model used:} Claude Sonnet (claude.ai, free tier)\\
\textbf{Date of debate:} March 2026\\

All responses were generated using the free tier interface at claude.ai. No system
prompt modifications were used. Each argument was presented in a new conversation
to prevent contextual carry-over between responses.

The objective of this exercise was not simply to obtain answers but to examine the
logical structure of the reasoning produced by the model. Where responses contained
logical weaknesses or incomplete reasoning these were identified explicitly and
analysed.

\section*{Overview}

This chapter records a structured debate with Claude across the four ethical conflict
points identified in Chapter~3. Each argument was prompted separately and Claude was
asked to provide at least two positions together with reasoning.

The responses are reproduced and then examined critically. Logical fallacies are
identified where they occur and the stronger position is selected based on evidence
from the literature and the ethical framework developed earlier in this report.

%-----------------------------------------------------------------------
\subsection*{Argument 1: Sensitivity Optimisation vs.\ Alert Fatigue}
\textit{Should the Threshold Be Raised?}

\noindent\textit{[Insert screenshot: Argument 1 prompt and response]}

\paragraph{Ethical Principle Evaluated.}~\\

\textbf{School:} Utilitarianism (Mill, principle of harm minimisation)

\textbf{Constraint:} The burden of false positive alerts must remain within
operational tolerance so that clinicians remain responsive to genuine warnings.
Specifically, if the clinician override rate exceeds 70 percent across consecutive
alert sessions, the system is no longer producing net benefit (Goal U2, Chapter~2).

\textbf{Violation condition:} If false positive alerts become sufficiently frequent
that clinicians ignore them reflexively, the system undermines patient safety rather
than supporting it.

\paragraph{Analysis.}

Recommendation 2 is the stronger position. Decisions about threshold adjustment
should be grounded in empirical evidence about clinician response behaviour.
Shwedeh and Alzoubi \cite{shwedeh2025} demonstrate that over-reliance and alert
fatigue emerge specifically in high-pressure deployment environments where governance
is insufficient, reinforcing that threshold decisions must be informed by local usage
data rather than modelling assumptions alone.

Recommendation 1 contains a logical flaw. The proposal assumes that ward type
reliably reflects the consequences of false positives and false negatives. A patient
with a risk score of six on a general ward has the same physiological condition as
a patient with the same score in an intensive care unit. Under this proposal only
one of those patients would generate an alert. General wards also typically have
less continuous monitoring than intensive care units, which may make missed
deterioration more dangerous rather than less. The logic of proportionality runs
in the opposite direction to what Claude claims.
\par\medskip
\textbf{Fallacy identified: Hasty Generalisation.} Ward location is treated as a
reliable proxy for patient risk when the relationship is not justified. Because the
proportionality claim depends on that assumption the argument collapses.
Recommendation 2 is therefore the stronger ethical position.

%-----------------------------------------------------------------------
\subsection*{Argument 2: Transparency vs.\ Proprietary Opacity}
\textit{Is the Responsibility Gap Defensible?}


\noindent\textit{[Insert screenshot: Argument 2 prompt and response]}

\paragraph{Ethical Principle Evaluated.}~\\

\textbf{School:} Deontological ethics (Kant, respect for rational agency)

\textbf{Constraint:} A clinician who is held legally and morally responsible for a
clinical decision must have meaningful epistemic access to the basis of the
information that prompted that decision. This is formalised as Goal D1 in
Chapter~2. The system satisfies this principle only if, for every alert fired, the
clinician can access at minimum the top contributing feature variables and their
directional influence on the score.

\textbf{Violation condition:} Any alert fired without accompanying feature-level
explanation constitutes a violation. In the current ESM deployment this condition
is violated for 100 percent of alerts.

\paragraph{Analysis.}

Position 2 identifies the central issue correctly. Clinicians remain accountable
for outcomes yet are denied access to the reasoning behind the alerts they must
interpret. Amann et al.\ \cite{amann2022} demonstrate that lack of explainability
causes systems nominally classified as decision support to drift toward decision
determining behaviour, undermining genuine human oversight. \v{C}artolovni et al.\
\cite{cartoloni2022} identify the resulting accountability gap as the dominant
ethical concern in the AI-based medical decision support literature.

Position 1 depends on an analogy with laboratory tests. This comparison is flawed.
A troponin value is a direct biological measurement with known, traceable error
properties. An AI-generated risk score is a learned inference derived from
correlations across many variables and may encode biases or data artefacts invisible
to the clinician \cite{cull2023}. Xu et al.\ \cite{xu2023} establish that post-hoc
explanations and population-level metrics are insufficient to guarantee
accountability in sub-symbolic clinical decision support systems. When the ESM is
systematically wrong for a particular patient group, that error is invisible at the
point of care.
\par\medskip
\textbf{Fallacy identified: False Analogy.} The analogy transfers the epistemic
legitimacy of laboratory measurements onto a fundamentally different class of
predictive system. Because the analogy is the foundation of Position~1 the defence
of opacity collapses with it. Position~2 is therefore the stronger position.

%-----------------------------------------------------------------------
\subsection*{Argument 3: Fairness Across Age Subgroups}
\textit{Should Testing Be Mandatory Pre-Deployment?}

\noindent\textit{[Insert screenshot: Argument 3 prompt and response]}

\paragraph{Ethical Principle Evaluated.}~\\

\textbf{School:} Deontological ethics (principle of equal treatment) and Utilitarian
justice

\textbf{Constraint:} Patients with equivalent clinical states must receive equivalent
diagnostic evaluation regardless of demographic subgroup. This is formalised in
Chapter~2 and operationalised in the Z3 fairness constraint as:
$\mathit{EquivalentClinicalState} \rightarrow (\mathit{Alert}_{elderly} =
\mathit{Alert}_{young})$. If this condition is violated it breaches the principle
of equal treatment under Goal D1.

\textbf{Violation condition:} Any systematic difference in alert rates between
clinically equivalent patients in different age subgroups constitutes a violation
detectable before deployment.

\paragraph{Analysis.}

Action 1 reflects the justice principle directly. Deploying without subgroup testing
accepts unknown differential harm. Elderly patients are not an edge case in sepsis
care and the physiological reasons their scores may be suppressed are well documented
clinically \cite{cull2023}.

Action 2 relies on an incorrect comparison class. The alternative to a
fairness-tested model is not the absence of clinical decision support. Most hospitals
already operate structured sepsis screening protocols and early warning scores.
Vasey et al.\ \cite{vasey2022} explicitly caution that ethical evaluation must
compare against the existing standard of care rather than against no intervention.

A second problem is that post-deployment monitoring collects evidence from patients
already being treated under the untested system. \v{C}artolovni et al.\
\cite{cartoloni2022} identify this as the accountability gap in AI-based medical
decision support, noting that deferring fairness assessment to post-deployment
monitoring makes institutional responsibility permanently ambiguous during the
observation period.
\par\medskip
\textbf{Fallacy identified: False Cause.} Action~2 implies that fairness violations
can only be discovered after deployment. The formal Z3 verification layer in this
project demonstrates that subgroup parity constraints can be logically checked before
any patient is involved. The UNSAT core output directly refutes the causal claim
that post-deployment monitoring is the only feasible approach. Action~1 is therefore
the stronger ethical position.

%-----------------------------------------------------------------------
\subsection*{Argument 4: Automation Bias}
\textit{Design Problem, Training Problem, or Inherent Risk?}

\noindent\textit{[Insert screenshot: Argument 4 prompt and response]}

\paragraph{Ethical Principle Evaluated.}~\\

\textbf{School:} Virtue ethics (Aristotle, phronesis and professional competence)

\textbf{Constraint:} Clinical decision support must assist rather than replace
independent clinical reasoning. This is formalised as Goals V1 and V2 in
Chapter~2. Automation bias is operationally present when clinician compliance with
alerts exceeds 90 percent without documented independent assessment. Deskilling
risk is flagged when independent override decisions fall below a minimum threshold
per clinician per month.

\textbf{Violation condition:} Compliance rate greater than 90 percent without
independent documentation, or a sustained reduction in clinician independent
override frequency over time.

\paragraph{Analysis.}

The distinction between interface design and training is reasonable. Panigutti
et al.\ \cite{panigutti2022} provide empirical evidence that explanations and
interface changes do influence reliance on AI advice, though not always in ways
that improve decision accuracy. Design and training therefore address different
failure modes and both are necessary.

However neither position engages with the deeper virtue ethics concern. Automation
bias is treated throughout as a cognitive shortcut problem. \v{C}artolovni et al.\
\cite{cartoloni2022} identify deskilling and loss of epistemic authority as
physician-specific risks that are distinct from momentary automation bias. A
clinician who has repeatedly followed ESM alerts without practising independent
diagnostic assessment may over time lose the calibrated clinical judgment that makes
meaningful override possible. This is Goal V2 in Chapter~2 and neither of Claude's
positions addresses it.

The closing observation that automation bias may be inherent is the most ethically
significant point in the response. Claude identifies the problem and then offers no
conclusion from it.
\par\medskip
\textbf{Fallacy identified: Argument from Ignorance.} The absence of a definitive
solution to inherent automation bias is implicitly treated as justification for not
requiring one. Both positions therefore remain incomplete because neither addresses
the long-term professional deskilling risk that virtue ethics identifies as the
deeper problem.

%-----------------------------------------------------------------------
\subsection*{Summary of Identified Errors}

\begin{table}[H]
\centering
\renewcommand{\arraystretch}{1.3}
\begin{tabular}{@{}p{2.5cm}p{3cm}p{3cm}p{4.5cm}@{}}
\toprule
\textbf{Argument} & \textbf{Claude Position} & \textbf{Fallacy} & \textbf{Supported Position} \\
\midrule
Threshold decision
  & Ward based threshold stratification
  & Hasty Generalisation
  & Empirical audit before any threshold change \\[4pt]
Model opacity
  & Analogy with laboratory tests
  & False Analogy
  & Explainability required when responsibility remains with the clinician \\[4pt]
Age fairness
  & Conditional deployment with monitoring
  & False Cause
  & Mandatory pre-deployment subgroup fairness testing \\[4pt]
Automation bias
  & Design and training sufficient
  & Argument from Ignorance
  & Additional safeguards required including explicit deskilling monitoring \\
\bottomrule
\end{tabular}
\caption{Summary of logical errors identified in Claude's responses.}
\label{tab:claude_errors}
\end{table}

Chapters~5 and~6 each take one of these fallacies and argue against Claude's
position using evidence from the repository.
%=======================================================================
\chapter{Teammate \#1 Says\ldots}
%=======================================================================

\section*{Selected Argument}

This chapter argues against Claude's Position~1 in Argument~2, which claimed that
opacity in the Epic Sepsis Model is ethically defensible because clinicians
interpreting an AI risk score are in the same position as clinicians interpreting
a laboratory result. This argument rests on a false analogy fallacy. Because the
fallacy undermines the entire defence of opacity, Position~2 that opacity is
ethically unacceptable when accountability remains with the clinician is the correct
position.

%-----------------------------------------------------------------------
\subsection*{Ethical Principle}

\textbf{School:} Deontological ethics (accountability and responsibility)

\textbf{Principle:}
A clinician who is held legally and morally responsible for a clinical decision must
have meaningful epistemic access to the basis of the information that informed that
decision.

\textbf{Operational constraint:}
If an AI system generates a clinical alert, the clinician must be able to access
the reasoning or contributing factors that produced the alert.

\textbf{Violation condition:}
If the system produces alerts without any explanation or accessible reasoning, the
clinician's responsibility cannot be justified. Under these conditions the system
creates a responsibility gap between formal authority and practical epistemic
control.

%-----------------------------------------------------------------------
\subsection*{The Fallacy: False Analogy}

Claude's Position~1 draws a structural equivalence between interpreting an ESM
alert and interpreting a troponin result. The argument is that in both cases the
clinician acts on a tool output without understanding its internal mechanics, and
that this is considered acceptable practice. The analogy fails because the two
situations have categorically different error structures.

A troponin result is a direct measurement of a biological quantity with documented,
investigable failure modes: reagent degradation, sample haemolysis, assay
interference from skeletal muscle troponin. When a troponin is anomalous or
inconsistent with clinical presentation, the clinician and the laboratory can reason
about why. The source of error is traceable and, in principle, correctable.

The ESM score is a learned inference from a proprietary model trained on
approximately 500,000 historical encounters, combining roughly 80 variables through
undisclosed statistical relationships \cite{cull2023}. Those relationships may
encode documentation biases, historical care inequities, or population-specific
confounders. When the ESM is systematically wrong for a particular patient group,
that error is invisible at the point of care. Xu et al.\ \cite{xu2023} establish
that this opacity is not a minor limitation but a fundamental ethical problem:
post-hoc explanations and population-level metrics are insufficient to guarantee
accountability in sub-symbolic CDSS. There is no equivalent of checking the reagent
batch. The clinician receives a number with no accessible basis for disagreement.

Equating these two situations is a false analogy. One error is traceable. The other
is structurally inaccessible.

%-----------------------------------------------------------------------
\subsection*{What the Fallacy Conceals}

The practical consequence of accepting Claude's false analogy is that the
responsibility gap is rendered invisible. Amann et al.\ \cite{amann2022} demonstrate
that lack of explainability causes systems nominally classified as decision support
to drift toward decision determining the clinician retains formal authority while
practical epistemic control has shifted to the model. Claude's analogy imports the
legitimacy of an established diagnostic measurement onto a system with fundamentally
different accountability implications, concealing this drift.

Consider the scenario the ESM is most likely to produce: a clinician receives an
alert for an elderly patient with atypical presentation, assesses the patient
independently, and decides the score does not reflect the clinical picture. They
document their reasoning and do not escalate. The patient later deteriorates. The
clinician's defence requires articulating not only what they observed but why they
disagreed with the score. Under an opaque system, the second part of that defence
has no grounding. This is not analogous to acting on a troponin without knowing the
biochemistry. It is acting on an inference without knowing what was inferred or why
precisely the condition \v{C}artolovni et al.\ \cite{cartoloni2022} identify as
the dominant accountability gap in AI-based medical decision support.

%-----------------------------------------------------------------------
\subsection*{The Structural Nature of the Gap}

The responsibility gap is therefore not an exceptional scenario but a structural
feature of the deployment architecture, not an edge case. The E4 constraint
formalises this directly:

\begin{lstlisting}[style=mystyle, language=Python]
def responsibility_gap(alert, explanation_available):
    if alert and not explanation_available:
        return True
    return False
\end{lstlisting}

In current published descriptions of the ESM deployment, clinicians are not provided
with feature-level explanations for alerts, meaning \texttt{explanation\_available}
is effectively \texttt{False} in practice. The function therefore returns
\texttt{True} whenever an alert is triggered. Claude argues that population-level
metrics such as 86\% sensitivity and 33.8\% PPV provide sufficient calibration.
They do not. Population statistics describe aggregate model behaviour, they say
nothing about the basis for any individual alert, which is what a clinician needs
to exercise the independent judgement that accountability requires.

%-----------------------------------------------------------------------
\subsection*{An Internal Contradiction}

Claude's Position~1 also contradicts what Claude acknowledged in Argument~4. The
defence of opacity depends on clinicians consciously using population-level metrics
to calibrate their response to each alert. But in Argument~4, Claude accepted that
clinicians under time pressure and fatigue default to following alerts without
independent assessment. Panigutti et al.\ \cite{panigutti2022} provide empirical
confirmation: explanations systematically increase human reliance on AI advice
regardless of whether reliance improves accuracy, meaning the calibrated reflective
clinician that Position~1 assumes is not the clinician that real deployment produces.
These two pictures of clinical behaviour are incompatible. Claude cannot use the
first to justify opacity while accepting the second as the operational reality.

%-----------------------------------------------------------------------
\subsection*{Conclusion}

The lab test analogy is a false analogy fallacy. It treats traceability of error as
irrelevant to accountability when it is central to it. Because the fallacy is the
load-bearing element of Position~1, the entire defence of opacity collapses with it.
Position~2 is correct: opacity is ethically unacceptable when clinicians bear full
accountability for decisions whose basis they cannot access. The responsibility gap
is not a theoretical concern, it is a structural feature of every alert the
system fires.

%=======================================================================
\chapter{Teammate \#2 Says\ldots}
%=======================================================================

\section*{Selected Argument}

This chapter argues against Claude's Non-action~2 in Argument~3, which proposed
replacing mandatory pre-deployment fairness testing with conditional deployment
accompanied by post-deployment monitoring. This argument contains a false cause
fallacy and it attributes the impossibility of pre-deployment fairness testing to
an epistemic limitation that does not exist. Because the fallacy is the basis for
the entire conditional deployment position, Action~1 mandating subgroup fairness
testing as a pre-deployment requirement is the correct position.

%-----------------------------------------------------------------------
\subsection*{Ethical Principle}

\textbf{School:} Deontological ethics (principle of equal treatment) with
utilitarian justice considerations.

\textbf{Principle:}
Patients with equivalent clinical states must receive equivalent diagnostic
evaluation regardless of demographic subgroup.

\textbf{Operational constraint:}
If two patients exhibit equivalent clinical indicators for sepsis, the system must
produce equivalent alert decisions regardless of age group.

\textbf{Violation condition:}
If two clinically equivalent patients receive different alert outcomes solely due
to subgroup membership, the fairness constraint is violated.

%-----------------------------------------------------------------------
\subsection*{The Fallacy: False Cause}

Claude's Non-action~2 implies that the reason for deferring fairness testing to
the post-deployment phase is that the violation cannot be known in advance. The
argument runs: a model with imperfect subgroup performance may still reduce net
sepsis mortality compared to no decision support, so conditional deployment with
post-deployment surveillance is a proportionate response to an uncertainty that
cannot be resolved before going live.

This is a false cause fallacy. The claimed cause of deferral epistemic impossibility
is not the actual cause. Fairness violations are formally detectable before
deployment. The Z3 verification layer in this project demonstrates this directly:

\begin{lstlisting}[style=mystyle, language=Python]
def unsat_fairness_violation():
    s = z3.Solver()
    Equivalent_Group = z3.Bool("Equivalent_Group")
    elderly_alert = z3.Bool("elderly_alert")
    young_alert = z3.Bool("young_alert")

    s.assert_and_track(
        z3.Implies(Equivalent_Group, elderly_alert == young_alert),
        "Fairness_Symmetry"
    )
    s.assert_and_track(Equivalent_Group, "Equivalent_Group_True")
    s.assert_and_track(elderly_alert, "Elderly_Alert_True")
    s.assert_and_track(z3.Not(young_alert), "Young_Alert_False")

    result = s.check()
    if result == z3.unsat:
        print("Unsat Core:", s.unsat_core())
\end{lstlisting}

This returns UNSAT with core \texttt{[Elderly\_Alert\_True, Equivalent\_Group\_True,
Fairness\_Symmetry, Young\_Alert\_False]}. The violation scenario is formally proven
to be logically inconsistent with the fairness constraint before any patient is
involved. Xu et al.\ \cite{xu2023} further identify the absence of formal
verification as a specific research gap in AI-based CDSS ethics, noting that ethical
compliance is seldom enforced through logic-based constraints which is precisely
what this implementation provides. The epistemic precondition for Claude's
conditional deployment argument is false.

%-----------------------------------------------------------------------
\subsection*{The Wrong Counterfactual}

Claude's argument also rests on a second error. The claim that a model with
imperfect subgroup performance may still reduce net sepsis mortality compared to
no decision support at all uses the wrong comparison class. The alternative to a
fairness-tested model is not an empty clinical environment. Most hospital systems
deploying the ESM already operate structured sepsis screening protocols, early
warning scores, and nursing assessment frameworks \cite{shwedeh2025}. Vasey et al.\
\cite{vasey2022} (DECIDE-AI) explicitly caution that high retrospective accuracy
does not guarantee clinical benefit or safety in real-world deployment, and that
ethical evaluation must account for the existing standard of care rather than a
zero baseline. The net benefit calculation Claude implies only holds if the
counterfactual is nothing, which it is not.

%-----------------------------------------------------------------------
\subsection*{Post-Deployment Monitoring Is Not a Safeguard}

Having removed the epistemic justification for deferral and corrected the
counterfactual, what remains of Non-action~2 is a governance process that collects
evidence of harm from patients who are already being managed under the untested
system. \v{C}artolovni et al.\ \cite{cartoloni2022} identify this accountability
gap directly: responsibility for harm in AI-based decision support is often unclear
across clinicians, developers, and institutions, and frameworks that defer fairness
assessment to post-deployment monitoring effectively make that ambiguity permanent
during the observation period. The patients harmed during monitoring are not data
points in a quality improvement exercise. They are people who did not receive a
timely alert because the institution chose not to test for the disparity before
deployment.

The justice principle identified under Goal D1 in Chapter~2 requires that
foreseeable disparate impacts be identified before they occur, not after. Elderly
patients present particular physiological challenges like blunted fever responses,
atypical inflammatory markers, baseline comorbidities that may systematically
suppress ESM scores \cite{cull2023}. This is not an unpredictable edge case. It is
a foreseeable risk that the pre-deployment formal verification layer in this project
is specifically designed to detect.

%-----------------------------------------------------------------------
\subsection*{Conclusion}

Claude's Non-action~2 rests on a false cause fallacy: the claim that fairness
violations cannot be known before deployment. The Z3 verification layer in this
project formally refutes that claim. With the fallacy removed, the conditional
deployment position has no epistemic justification, uses the wrong counterfactual,
and treats foreseeable harm to elderly patients as an administrative cost. Action~1
is correct, subgroup fairness testing must be a pre-deployment requirement, not
a post-deployment observation.

%=======================================================================
\chapter{Conclusions}
%=======================================================================

This project examined the Epic Sepsis Model through three ethical frameworks:
utilitarianism, deontological ethics, and virtue ethics. Each framework illuminated
a different dimension of the system's ethical profile. Utilitarian analysis surfaced
the tension between sensitivity optimisation and false positive burden. Virtue ethics
raised the longer term concern about automation bias and clinical deskilling.
Deontological ethics identified the accountability and governance failures that sit
at the core of the system's design.

\medskip

Of the three, deontological ethics proved the most tractable for both expressing the
problems and implementing solutions. The reason is structural. Deontological
reasoning operates through obligations and rules: if an alert fires, it must be
acknowledged; if a clinician is accountable, they must have epistemic access; if two
patients are clinically equivalent, they must receive equivalent treatment. These
translate directly into executable constraints. Every ethical goal in the
implementation, from the responsibility gap check to the governance constraint to
the fairness symmetry rule, is a formalised duty of this kind. The Z3 verification
layer makes this particularly clear: encoding an ethical obligation as a logical
formula and checking whether it is satisfiable is deontological reasoning in a
formal system.
\medskip
Utilitarianism was useful for framing population level trade-offs, particularly
around threshold configuration and false positive rates, but harder to implement
rigorously. Meaningful utilitarian validation would require outcome data across real
patient populations, which a simulation abstraction cannot provide. The threshold
trade-off simulation approximates this but cannot quantify net mortality benefit,
which limits how far the utilitarian analysis can be taken within this scope.
\medskip
Virtue ethics raised the most important long term concern, that repeated AI assisted
workflows may erode the professional judgment that makes meaningful human oversight
possible, but was the hardest to encode. Deskilling is a dynamic process that
unfolds over years of clinical practice. It does not reduce naturally to a binary
constraint or a satisfiability check, which meant it remained at the level of
analysis rather than implementation.
\medskip
The broader lesson from this project is that ethical goals in AI systems are most
effectively addressed when they are treated as first class design constraints rather
than post-hoc commentary. The cases where the implementation fell short, particularly
around the responsibility gap and the fairness simulation, were precisely the cases
where the real system's opacity made formal constraint verification difficult. That
difficulty is itself an ethical signal.
%=======================================================================
\chapter*{Acknowledgements}
%=======================================================================

The authors would like to thank Dr.\ Vivek Nallur, School of Computer Science,
University College Dublin, for his supervision, guidance, and feedback throughout
this project. His direction on ethical framework construction, the debate structure,
and the formal verification approach was instrumental in shaping the analysis
presented in this report.

\medskip
The authors also wish to acknowledge the researchers whose published work made this
project possible. The validation study by Cull et al.\ \cite{cull2023} provided the
clinical foundation for the case study. The broader literature on interpretability,
explainability, governance, and ethical evaluation of clinical AI systems informed
the ethical framework developed in Chapter~2 and the critical analysis presented in
Chapters~4, 5, and~6.

\medskip
This project was submitted in part fulfilment of the MSc in Advanced AI at
University College Dublin, March 2026.
%=======================================================================

\printbibliography

\end{document}
